% Generated draft: requires mathematical and source review.
\documentclass[11pt]{article}
\usepackage[margin=25mm]{geometry}
\usepackage{fontspec}
\setmainfont{Latin Modern Roman}
\usepackage{amsmath,amssymb,mathtools,unicode-math}
\setmathfont{Latin Modern Math}
\usepackage{enumitem,needspace,xurl,hyperref}
\setlength{\parindent}{0pt}
\setlength{\parskip}{5pt}
\setlength{\emergencystretch}{3em}
\newcommand{\N}{\mathbb{N}}
\newcommand{\Z}{\mathbb{Z}}
\newcommand{\Q}{\mathbb{Q}}
\newcommand{\R}{\mathbb{R}}
\newcommand{\C}{\mathbb{C}}
\newcommand{\F}{\mathbb{F}}
\newcommand{\A}{\mathbb{A}}
\newcommand{\PP}{\mathbb{P}}
\newcommand{\E}{\mathbb{E}}
\newcommand{\eps}{\varepsilon}
\newcommand{\dd}{\,\mathrm{d}}
\begin{document}
\section*{UnsolvedMath ID: 30006597}
\textbf{Source problem code:} catalogue entry
\subsection{Definitions and mathematical statement}
Let \((a_n)_{n\geq 0}\) be a sequence of integers, and define formal power series
\[
f(x)=\sum_{n\geq 0}a_nx^n\in\Z[[x]],\qquad F(x)=\sum_{n\geq 0}\binom{2n}{n}a_nx^n,
\qquad \binom{2n}{n}=\frac{(2n)!}{(n!)^2}.
\]
Algebraicity of \(f\) means that there exists a nonzero polynomial \(P(x,y)\in\Q[x,y]\) such that \(P(x,f(x))=0\) as a formal power series.

The supplied question asks whether, for every such integer sequence, the existence of a second-order homogeneous linear differential equation with polynomial coefficients and regular singularities satisfied by \(F\) implies that \(f\) is algebraic. Such an equation has the form
\[
p_2(x)F''(x)+p_1(x)F'(x)+p_0(x)F(x)=0,
\qquad p_2\not\equiv 0,
\]
where derivatives and the identity may be interpreted formally. The source does not require this equation to have minimal order among equations annihilating \(F\).

\textbf{Unresolved differential-equation conventions.} The supplied record does not specify the coefficient field of the polynomials or explicitly state whether regularity is required at infinity. To express the remaining choices precisely, fix a subfield \(K\subseteq\C\), require \(p_i\in K[x]\), and write
\[
q(x)=\frac{p_1(x)}{p_2(x)},\qquad r(x)=\frac{p_0(x)}{p_2(x)}.
\]
At a finite point \(c\in\C\), the usual regular-singular condition is that \((x-c)q(x)\) and \((x-c)^2r(x)\) extend holomorphically to \(c\). Ordinary points satisfy this condition as well. If infinity is included, the additional condition is
\[
q(x)=O(x^{-1}),\qquad r(x)=O(x^{-2})\qquad (x\to\infty).
\]
For each fixed choice of \(K\) and of whether infinity is included, the question is therefore whether every \(f\in\Z[[x]]\) whose associated \(F\) satisfies an equation with these properties must be algebraic over \(\Q(x)\). These choices remain qualifications requiring source review; no particular choice is asserted here to be the intended one. No convergence assumption on either series is supplied.

\textbf{Conflicting original extraction.} The record's original-statement field also contains the request
\[
\text{if }k\geq 1\text{ and }29\nmid k,\quad b(29k+21)\equiv 0\pmod 4.
\]
The sequence \(b\) is undefined, and no relation to \(f\), \(F\), or \((a_n)\) is given. This additional request cannot be formulated completely from the supplied evidence and is not incorporated as a hypothesis or conclusion of the central-binomial question.
\subsection{Short English statement}
Start with an integer-coefficient formal power series and multiply its \(n\)-th coefficient by \(\binom{2n}{n}\). The supplied question asks whether a second-order differential equation with polynomial coefficients and regular singularities for the resulting series forces the original series to be algebraic. The coefficient field, the treatment of infinity, and an unrelated congruence involving an undefined sequence require clarification from the sources. The supplied record labels the question open; that status is not independently verified here.
\subsection{Sources}
\begin{itemize}
\item Local registry: \texttt{data/unsolvedmath/problems.json}, permanent UnsolvedMath ID \(30006597\), title \emph{Algebraicity Under Central-Binomial Untwisting}, exact source problem code \texttt{catalogue entry}. No individual proposer is identified in the supplied record.
\item Supplied report citation: \emph{Multiple Zeta Values and Modular Forms} (2026), Oberwolfach Reports, Report 7/2026, pp.\ 433--486, DOI \texttt{10.4171/owr/2026/7}. \url{https://doi.org/10.4171/owr/2026/7}. The supplied background describes Oberwolfach Reports as documenting mathematical workshops and their open-problem discussions.
\item Supplied literature reference: van Hoeij--van Straten--Zudilin, \emph{Hyperelliptic saga}. \url{https://doi.org/10.48550/arXiv.2306.04921}.
\item Supplied literature reference: Journal of Experimental Mathematics version. \url{https://jexpmath.org/index.php/jem/article/view/Vol-1Issue-2Paper-5}.
\item Supplied dated curation, checked exactly \texttt{2026-08-22}: the record labels the problem open and states that the 2025 primary paper formulates the same question, reports evidence against a broad positive answer, and supplies neither a proof nor a counterexample. This is the supplied dated assessment, not an independently verified literature conclusion. The background explicitly calls for independent verification of time-sensitive status and cited claims.
\item Supplied statement-verification metadata: \texttt{corrected\_verified}; the record states that the question was reconstructed and checked against the official Oberwolfach report. This is a supplied verification claim, not a source check performed for this draft.
\end{itemize}
\end{document}
